% 07_evidence.tex — Empirical Evidence
% ============================================================

\section{Empirical Evidence}
\label{sec:evidence}

We present three tiers of empirical validation: a 26-modality breadth
benchmark, a 16-modality operator-correction study, and a CASSI deep
dive with per-solver, per-scenario analysis.

% ============================================================
%  7.1  26-Modality Benchmark
% ============================================================
\subsection{26-Modality Benchmark}
\label{subsec:benchmark}

\Cref{tab:26mod} reports the peak signal-to-noise ratio (PSNR) achieved
by the \pwm{} pipeline across all 26~registered modalities under
Scenario~I (ideal operator, no mismatch).  Of the 64 modalities registered in the OperatorGraph IR, 26 have been numerically benchmarked to date.  Every modality passes the
benchmark threshold; no modality is excluded or deferred.

\begin{table}[t]
\centering
\caption{26-modality benchmark results (Scenario~I, ideal operator).
All 26/26 modalities pass.  Average PSNR $\approx 37.0$~dB (excluding Phase Retrieval identity test); range 25.46--64.84~dB over physically meaningful modalities.}
\label{tab:26mod}
\smallskip
\begin{tabular}{@{}rlr@{\qquad}rlr@{}}
\toprule
\textbf{\#} & \textbf{Modality} & \textbf{PSNR (dB)}
&
\textbf{\#} & \textbf{Modality} & \textbf{PSNR (dB)} \\
\midrule
 1 & Widefield              & 27.31 & 14 & Holography               & 46.85 \\
 2 & Widefield Low-Dose     & 32.88 & 15 & NeRF~\cite{mildenhall2020nerf} & 61.35 \\
 3 & Confocal Live-Cell     & 29.80 & 16 & 3D Gaussian Splatting~\cite{kerbl2023gaussiansplatting} & 30.89 \\
 4 & Confocal 3D            & 29.01 & 17 & Matrix (Generic)         & 33.86 \\
 5 & SIM                    & 27.48 & 18 & Panorama Multifocal      & 27.90 \\
 6 & CASSI~\cite{wagadarikar2008cassi} & 34.81 & 19 & Light Field              & 30.35 \\
 7 & SPC~\cite{duarte2008spc}       & 28.86 & 20 & Integral Photography     & 27.85 \\
 8 & CACTI~\cite{llull2013cacti}     & 35.33 & 21 & Phase Retrieval          &\best{100.00}\\
 9 & Lensless               & 26.85 & 22 & FLIM                     & 35.38 \\
10 & Light-Sheet            & 26.05 & 23 & Photoacoustic            & 50.54 \\
11 & CT                     &\worst{25.46}& 24 & OCT                  & 64.84 \\
12 & MRI~\cite{zbontar2018fastmri}   & 44.97 & 25 & FPM                      & 34.57 \\
13 & Ptychography~\cite{maiden2009ptychography} & 59.41 & 26 & DOT  & 32.06 \\
\bottomrule
\end{tabular}
\end{table}

The range spans nearly 75~dB, reflecting the diversity of forward models
--- from the heavily ill-posed computed tomography (CT~\cite{chen2017redcnn},
25.46~dB) to phase retrieval, where the analytically invertible
forward model yields an identity reconstruction
(100.00~dB, a numerical ceiling rather than a measured value).  The average of ${\approx}\,37.0$~dB (excluding the Phase Retrieval identity test) confirms that the
unified pipeline delivers competitive reconstruction quality across
fundamentally different physics.

% ============================================================
%  7.2  16-Modality Operator Correction
% ============================================================
\subsection{16-Modality Operator Correction}
\label{subsec:correction}

\Cref{tab:correction} evaluates the \pwm{} operator-correction module
across 16~modalities.  For each modality we report the PSNR
\emph{before} correction (Scenario~II, mismatched operator) and
\emph{after} correction (Scenario~III, corrected operator), together
with the absolute improvement.

\begin{table}[t]
\centering
\caption{Operator correction across 16~modalities.  ``Before'' is
Scenario~II (mismatched operator); ``After'' is Scenario~III (corrected
operator).  Improvements range from $+3.55$ to $+48.25$~dB.  Phase~2
and Phase~4 entries are registered but not yet numerically evaluated.}
\label{tab:correction}
\smallskip
\begin{tabular}{@{}llccc@{}}
\toprule
\textbf{Modality} & \textbf{Mismatch Type}
    & \textbf{Before (dB)} & \textbf{After (dB)} & \textbf{$\Delta$ (dB)} \\
\midrule
Matrix (Generic)     & gain\_bias             & 11.14 & 23.35 & $+$12.21 \\
CT                   & center\_of\_rotation   & 13.41 & 24.09 & $+$10.67 \\
CACTI~\cite{llull2013cacti,wang2023efficientsci} & mask\_timing & 14.48 & 37.42 & \best{$+$22.94} \\
Lensless             & psf\_shift             & 23.48 & 27.03 & \worst{$+$3.55}  \\
MRI                  & coil\_sensitivities    &  6.94 & 55.19 & \best{$+$48.25} \\
SPC$^\dagger$ & gain\_bias   & 11.14 & 23.35 & $+$12.21 \\
CASSI (Alg~1)        & mask\_geo+dispersion   & 15.79 & 21.20 & $+$5.40  \\
CASSI (Alg~2)        & mask\_geo+dispersion   & 15.79 & 21.54 & $+$5.74  \\
Ptychography         & position\_offset       & 17.35 & 24.44 & $+$7.09  \\
\midrule
\multicolumn{5}{@{}l}{\textit{Phase~2 additions}} \\
OCT                  & dispersion             & \multicolumn{3}{c}{registered} \\
Light Field          & depth\_estimation      & \multicolumn{3}{c}{registered} \\
\midrule
\multicolumn{5}{@{}l}{\textit{Phase~4 additions}} \\
DOT                  & scatter\_coeff         & \multicolumn{3}{c}{registered} \\
Photoacoustic        & speed\_of\_sound       & \multicolumn{3}{c}{registered} \\
FLIM                 & irf\_shift             & \multicolumn{3}{c}{registered} \\
Integral Photography & disparity\_offset      & \multicolumn{3}{c}{registered} \\
FPM                  & led\_position          & \multicolumn{3}{c}{registered} \\
\bottomrule
\end{tabular}
\end{table}

{\footnotesize $\dagger$~SPC uses the same gain-bias operator template
as the generic Matrix modality, yielding identical correction
performance.}

\medskip
\noindent
\textbf{Note on CASSI correction algorithms.}
In \cref{tab:correction}, Alg~1 refers to hierarchical beam search (coarse correction) and Alg~2 refers to joint gradient refinement (fine correction); see Methods in the companion paper~\cite{yang2026inversenet} for details.

\medskip
\noindent
\textbf{Note on CASSI baselines.}
The Scenario~II baseline for CASSI in \cref{tab:correction} (15.79~dB) is lower than the corresponding value in the sensitivity table (\cref{tab:sensitivity}, 18.09~dB for MST-L) because the correction study uses a multi-parameter mismatch (mask geometry + dispersion), which is more severe than the 1-pixel-shift-only mismatch used in the sensitivity analysis.

\medskip
\noindent
The most dramatic improvement is MRI coil-sensitivity correction at
$+48.25$~dB.  This large gain reflects a severe mismatch scenario
(Scenario~II at 6.94~dB is near the noise floor); the corrected
55.19~dB exceeds published baselines (42.3~dB) because the mismatch
was synthetically injected on a high-quality forward model.
Even the smallest gain (Lensless PSF shift, $+3.55$~dB) represents a
visually significant improvement.  These results confirm that operator correction is not a
marginal refinement but a \emph{first-order} determinant of
reconstruction quality.

% ============================================================
%  7.3  CASSI Deep Dive
% ============================================================
\subsection{CASSI Deep Dive}
\label{subsec:cassi}

We conduct a detailed study of coded aperture snapshot spectral imaging
(CASSI) to illustrate how the Triad Law (\cref{sec:triad}) manifests in
a single modality.  \Cref{tab:cassi} reports PSNR for five solvers
across three scenarios plus a projected fine-tuned estimate:
%
\begin{itemize}
    \item \textbf{Scenario~I} --- ideal (true) operator;
    \item \textbf{Scenario~II} --- mismatched (nominal) operator;
    \item \textbf{Scenario~III} --- \pwm{}-corrected operator;
    \item \textbf{Estimated} --- corrected operator with additional solver
          fine-tuning (projected, not a formal scenario).
\end{itemize}
%
\noindent Scenario~IV (Oracle) reproduces Scenario~I by design for CASSI and is omitted for brevity.

\begin{table}[t]
\centering
\caption{CASSI deep dive: 5~solvers $\times$ 3~scenarios + projected fine-tuned estimate.  The mismatch
cost (I$\to$II) dominates the solver upgrade gain under ideal conditions.  Scenario~IV (Oracle) reproduces Scenario~I by design and is omitted for brevity.}
\label{tab:cassi}
\smallskip
\begin{tabular}{@{}lcccc@{}}
\toprule
\textbf{Solver}
    & \textbf{I (Ideal)}
    & \textbf{II (Mismatch)}
    & \textbf{III (Corrected)}
    & \textbf{Est.\ Fine-tuned}$^\dagger$ \\
\midrule
GAP-TV~\cite{yuan2016gaptv}   & 24.22 & 19.66 & 20.17 & ${\approx}$23.0 \\
PnP-HSICNN~\cite{wang2017adaptive} & 25.12 & 19.10 & 19.81 & ${\approx}$24.0 \\
MST-S~\cite{cai2022mst}     & 33.98 & 18.01 & 21.01 & ${\approx}$30.0 \\
MST-L~\cite{cai2022mst}     & 34.81 & 18.09 & 21.10 & ${\approx}$31.0 \\
HDNet~\cite{hu2022hdnet}     & 34.66 & 24.18 & 24.23 & ${\approx}$33.0 \\
\bottomrule
\end{tabular}
\end{table}

\smallskip
{\footnotesize $\dagger$~Projected values based on expected fine-tuning gains; not measured.}

\paragraph{Key insight.}
The mismatch cost --- the gap from Scenario~II to Scenario~I --- is
catastrophic.  For MST-L the drop is $34.81 - 18.09 = 16.72$~dB, while
the best solver upgrade under ideal conditions (GAP-TV $\to$ MST-L)
yields $34.81 - 24.22 = 10.59$~dB.  The mismatch penalty is
\textbf{$1.58\times$} the solver upgrade gain.  Stated differently:
\emph{fixing the operator is worth more than upgrading from a classical
solver to the state of the art}.

Moreover, the learned solvers (MST-S, MST-L) are \emph{more} sensitive
to mismatch than the model-based solvers (GAP-TV, PnP-HSICNN).  MST-L
drops by 16.72~dB versus GAP-TV's 4.56~dB.  This is consistent with the
hypothesis that data-driven methods overfit to the training-time operator
and generalize poorly when the physical system drifts.

% ============================================================
%  SolveEverything gears
% ============================================================
\subsection{SolveEverything Implementation Status}

In addition to the 9-layer Industrial Intelligence Stack (\cref{sec:stack}), the SolveEverything framework~\cite{wissner2026solveeverything} defines 10 operational gears---concrete mechanisms that drive each layer.  \Cref{tab:gears} reports the implementation status of each gear within PWM.

\begin{table}[t]
\centering
\caption{SolveEverything 10-gear implementation status.}
\label{tab:gears}
\smallskip
\begin{tabular}{@{}clll@{}}
\toprule
\textbf{Gear} & \textbf{Name} & \textbf{Status} \\
\midrule
 1 & Targeting System    & \textcolor{built}{\textbf{BUILT}}    \\
 2 & Outcome Contracts   & \textcolor{partial}{\textbf{PARTIAL}} \\
 3 & Compute Escrow      & \textcolor{partial}{\textbf{PARTIAL}} \\
 4 & Action Networks     & \textcolor{planned}{\textbf{PLANNED}} \\
 5 & Data Trusts         & \textcolor{built}{\textbf{FOUNDATION}} \\
 6 & Decision Logs       & \textcolor{built}{\textbf{BUILT}}    \\
 7 & Two-Source Rule     & \textcolor{partial}{\textbf{PARTIAL}} \\
 8 & Compute + Energy    & \textcolor{planned}{OUT OF SCOPE}     \\
 9 & Fairness Targets    & \textcolor{partial}{\textbf{PARTIAL}} \\
10 & Literacy            & \textcolor{partial}{\textbf{PARTIAL}} \\
\bottomrule
\end{tabular}
\end{table}

Two gears are fully built (Targeting System, Decision Logs) and one has its foundation laid (Data Trusts).  Five are partially implemented, one is
planned, and one (Compute + Energy) is currently out of scope for the
imaging-focused release.  The partial gears are under active development
and are expected to reach full status before the public v1.0 release.
